% ---------------------------------------------------------------------------
% Draft prose for the BAS baseline, written 2026-08-14 from the measured results.
%
% Two fragments: a Methods subsection (for sections/04_methods.tex) and a Results
% subsection (for sections/02_results.tex). Every number here is reproducible from
% scripts/bas/; the supporting measurements are recorded in docs/bas-findings.md.
%
% Handed over rather than inserted, per paper/HANDOFF_TO_CODING_AGENT.md. Note that the
% existing text at 02_results.tex:295 introduces the BAS baseline as "a video backbone
% extracting clip features across the full-match feed", which does not describe this
% system and would contradict these tables.
% ---------------------------------------------------------------------------


% ===========================================================================
% METHODS  ->  sections/04_methods.tex
% ===========================================================================

\subsection{Trajectory-based ball action spotting}
\label{subsec:methods_bas_traj}

The ball action spotting baseline predicts, for each of the twelve event classes, a ranked
list of timestamps at which an event of that class occurs. We report temporal mean average
precision at a tight tolerance of $\tau = 1\,\text{s}$ and a loose tolerance of
$\tau = 5\,\text{s}$, following the SoccerNet ball action spotting
protocol~\cite{giancola2018soccernet,deliege2021soccernet}: a prediction counts as a true
positive when an as-yet-unmatched ground-truth event of the same class lies within $\tau$ of
it in the same period. Because the class distribution spans a factor of $301$ between the
most and least frequent class, we report both the unweighted mean over the twelve classes and
a mean weighted by class support, and we give the support alongside every per-class figure.
Average precision is computed over predictions pooled across the test matches rather than
averaged per match, so that a rare class is scored on all of its instances at once.

\paragraph{Benchmark definition.}
Three of the ten matches (117092, 132831 and 132877) were played as three periods of
$45\,\text{min}$ rather than two, and no video or game-state annotation was produced for the
third period of any of them. The $2{,}231$ events that fall in those periods --- $9.4\%$ of
the $23{,}663$ annotated events, and $22.8\%$ of the annotations of test match 132831 alone
--- therefore have no corresponding input of any kind. Scoring against them depresses recall
by an amount that differs from match to match: a prediction that reproduces every event with
an observable counterpart scores $0.841$ rather than $1.000$ at $\tau = 1\,\text{s}$, with the
entire shortfall falling on 132831. We consequently define the benchmark over the two filmed
periods, comprising $21{,}432$ events, of which $17{,}060$ fall in the eight training and
validation matches and $4{,}372$ in the two test matches.

Assigning an event to a period requires care. The \texttt{gameTime} field carries a period
prefix, but that prefix is contradicted by the event's own timestamp for $102$ events, and the
timestamp alone cannot substitute for it because consecutive periods overlap on the nominal
match clock --- the first half of 118576 runs to $48{:}29$ while its second half begins at
$45{:}00$. We therefore treat the prefix as a hypothesis and accept it only when the resulting
frame index falls inside the annotated frame range of that period, which is precisely the
condition that the event has input data.

\paragraph{Temporal alignment.}
The mapping from an event timestamp to a game-state frame is
$f = 1 + (p - t_0)/40$, where $p$ is the annotated position in milliseconds and $t_0$ is the
match-clock time of the first annotated frame of that period. We determine $t_0$ by
measurement rather than from the recording metadata. For each candidate offset we compute the
fraction of events for which the annotated actor is the player closest to the ball, and take
the offset that maximises it; this statistic reaches $0.86$ to $0.95$ at its optimum against a
floor of $0.13$ to $0.25$, and recovers deliberately planted offsets exactly. The measured
offsets are bimodal: six halves require no correction and fourteen require exactly one second,
with no intermediate value, which matches the distribution of frame offsets independently
measured between the annotations and the released video. Deriving $t_0$ from the metadata
instead misaligns fourteen of the twenty halves by one second and costs $0.12$ mean average
precision at $\tau = 1\,\text{s}$.

\paragraph{Input representation.}
The model consumes player positions on the pitch and no pixels. Positions are quantised to
$1.05\,\text{m}$ in $x$ and $0.68\,\text{m}$ in $y$, so a single-frame finite difference
carries $26\,\text{m}\,\text{s}^{-1}$ of quantisation noise; velocities are therefore computed
at two scales, $0.48\,\text{s}$ and $1.52\,\text{s}$, and both are supplied to the model. From
each frame we derive $82$ features in four groups: the centroid, dispersion and convex hull of
all $22$ players; a soft minimum over opposing player pairs, which serves as a proxy for the
location of the ball; per-team centroid, dispersion, hull, speed and defensive and attacking
lines; and a $6 \times 3$ bilinearly-binned player-density grid per team, which encodes
configuration without depending on any player ordering. Features are sampled at $5\,\text{Hz}$,
comfortably finer than the tighter evaluation tolerance.

We report two input conditions. The \emph{trajectory} condition uses the $82$ features above
and is reproducible from the released game-state annotations alone. The \emph{trajectory and
ball} condition adds $19$ features derived from the provider's ball track --- position,
velocity, speed, acceleration, distance to the nearest player of each team, the number of
players within $3$, $5$ and $10\,\text{m}$, and distances to each goal and to the boundary
lines. The per-frame \texttt{ballStatus} field is deliberately excluded: its \texttt{BALLOUT}
value is close to a direct label for the \texttt{Out} class on eight matches, and it is
constant on the remaining two.

\paragraph{Model, training and decoding.}
A temporal convolutional network maps the feature sequence to twelve per-frame logits. It
comprises a pointwise input projection to $64$ channels, six residual blocks with dilations
$1$ to $32$ giving a receptive field of $\pm 25\,\text{s}$, and a pointwise output projection.
The twelve outputs are independent sigmoids rather than a softmax over classes, because
simultaneous events are part of the annotation convention. Targets are Gaussian of standard
deviation $300\,\text{ms}$ centred on each event, and the loss is binary cross-entropy with
per-class positive weighting capped at $50$. Training uses random $102\,\text{s}$ crops with
reflection augmentation in both pitch axes, which is label-preserving because no class is
handed. Models are selected on validation mean average precision rather than validation loss,
which the positive weighting renders non-monotone in the reported metric. At inference,
per-class local maxima above a floor are thinned by within-class non-maximum suppression, and
each surviving peak is emitted with its height as its confidence. The floor, the suppression
radius, the model capacity and the number of epochs are chosen on the validation matches of
the fold being evaluated, never on its test matches.

\paragraph{Split.}
We evaluate with five-fold cross-match cross-validation. Each fold holds out two matches for
test and two for validation, every match appears in the test set exactly once, and fold~0 is
the SoccerTrack Challenge 2025 pair so that it remains directly comparable with the
leaderboard. Validation for each fold is the following fold's test pair; the rule is arbitrary
but fixed, and no fold tunes on its own test matches. One consequence is that fold~4 validates
on the Challenge pair, so the cross-validated mean is not fully independent of the Challenge
fold --- a standard property of $k$-fold that we state rather than conceal.

We deliberately do not split within matches. The two halves of a match share the same
twenty-two players, both squads' shape, the pitch, the camera and the weather, and --- more
consequentially for this release --- every per-match data property, including the y-axis
convention, the ball clamping and the temporal offset described above. A within-match split
would place all of those on both sides of the partition and so conceal exactly the variation
that a held-out match exposes. We note that a fully team-disjoint split is not achievable with
these ten fixtures: 筑波大学-B appears in four matches and 筑波大学-C1 in three, so some folds
share a team between training and test.

\paragraph{Chance baseline.}
Because two classes account for $82\%$ of events, and a \texttt{Pass} occurs on average every
$2.4\,\text{s}$ of tracked play, evenly spaced guesses achieve a non-trivial score. We
therefore report a uniform-cadence baseline that emits spots at each class's mean rate
estimated on the training matches and knows nothing else about the match.


% ===========================================================================
% DATASET AND SPLIT  ->  sections/02_results.tex, before the results
% ===========================================================================

\paragraph{Dataset and split.}
The ten matches comprise $1{,}400{,}774$ tracked frames, or $934\,\text{min}$ of play at
$25\,\text{fps}$, and $30.8$ million player-frames at twenty-two tracked players per frame.
Table~\ref{tab:bas_dataset} gives the per-match breakdown. We adopt the SoccerTrack Challenge
2025 split, holding out 128057 and 132831 for test; of the remaining eight matches we reserve
117093 and 132877 for validation and train on the other six. The validation pair is chosen to
mirror the composition of the test pair, each containing one two-period and one three-period
match, so that every hyperparameter, the decoding threshold and the stopping epoch are chosen
under the same conditions the test split presents. The resulting benchmark contains $21{,}432$
events, of which $12{,}530$ are used for training, $4{,}530$ for validation and $4{,}372$ for
test.

Of the $23{,}663$ annotated events, $2{,}231$ fall in a third period that was played but
neither filmed nor tracked, and are excluded as described above. Actor linkage is high and
close to uniform across matches: $99.5\%$ of benchmark events name the player who performed
them, with a per-match range of $98.6$ to $99.9\%$, which is what permits an event to be
joined to that player's trajectory.

Two per-match limitations bear on the results and are marked in
Table~\ref{tab:bas_dataset}. The recorded ball track is clamped to the pitch rectangle in
132831 and 132877, so it never leaves the field of play in either; 132831 is in the test
split, and the consequence is quantified in Section~\ref{subsec:bas_results}. The panoramic
video is also not a single format across the release, varying between
$4096\times1080$, $4096\times1084$, $3840\times1906$ and $3840\times1504$.


% ===========================================================================
% RESULTS  ->  sections/02_results.tex
% ===========================================================================

\subsection{Ball action spotting from game state alone}
\label{subsec:bas_results}

Table~\ref{tab:bas_cv} reports the five-fold cross-match evaluation and
Table~\ref{tab:bas_per_class} the per-class breakdown pooled over all ten matches. All
figures are computed on ground-truth player trajectories, so the model assumes perfect
game-state reconstruction and is an upper bound on any system that must estimate the tracks
first. The spotter is also non-causal, seeing $\pm 25\,\text{s}$ around each prediction, which
is standard for offline action spotting but does not describe a real-time system.

From player positions alone the model attains a cross-validated $0.417 \pm 0.028$ mean average
precision at $\tau = 1\,\text{s}$ and $0.574 \pm 0.048$ at $\tau = 5\,\text{s}$, against
$0.025$ and $0.100$ for the uniform-cadence baseline. Adding the ball track raises these to
$0.662 \pm 0.065$ and $0.701 \pm 0.041$. The chance baseline is not a formality: its
support-weighted score at $\tau = 5\,\text{s}$ is $0.403$, so any weighted figure below that
is worse than guessing, and the margin rather than the absolute value is the quantity of
interest.

\paragraph{Why we cross-validate rather than report a single split.}
With ten matches, which pair is held out matters more than any modelling choice we make. The
standard deviation across the five folds is $0.028$ (trajectory) and $0.065$ (trajectory and
ball) at $\tau = 1\,\text{s}$, against $0.006$ and $0.028$ across three training seeds on a
fixed split --- a factor of $4.3$ and $2.3$ respectively. The Challenge fold is the weakest of
the five for the trajectory condition ($0.391$ against a mean of $0.417$). A single $8/2$ draw
reported without this context invites a fold effect to be read as a method effect.

The choice of \emph{validation} matches is equally consequential and less obvious. Holding the
test pair fixed at the Challenge matches and changing only the validation pair from
$\{117093, 132877\}$ to $\{117092, 117093\}$ moves the with-ball figure from $0.519$ to
$0.670$, five times the seed standard deviation. The cause is identifiable: 132877 is one of
the two matches whose recorded ball is clamped to the pitch rectangle, so a with-ball model
tuned on it selects a decoding threshold suited to censored ball data. We therefore fix the
validation rule as part of the protocol --- each fold validates on the following fold's test
pair --- and report the cross-validated mean rather than any single configuration.

\paragraph{Per-class structure.}
The aggregate hides the finding. Events defined by a duel outcome remain the two weakest
classes in both conditions: \texttt{Ball Player Block} reaches $0.117$ from trajectories and
$0.240$ with the ball, and \texttt{Player Successful Tackle} $0.088$ and $0.196$, against
$0.554$ and $0.868$ for \texttt{Pass}. Both ask which of two converging players won a contact,
and neither the players' ground positions nor the ball's location distinguishes the outcome.
This is a property of two-dimensional tracking rather than of the model, and it is the clearest
statement in these results of what the modality cannot represent.

The ball helps everywhere, but most where the event is defined by the ball itself rather than
by the players. \texttt{Throw In} gains $0.44$ and \texttt{Free Kick} $0.41$ --- restarts, at
which the ball's position and stillness are decisive and the players' configuration is
ambiguous --- while \texttt{Pass} and \texttt{Drive}, which the player configuration already
signals, gain $0.31$ and $0.30$ from a much higher base. \texttt{Out} remains comparatively
weak at $0.604$ even with the ball, which is consistent with two of the ten matches recording
a ball that never leaves the field of play.

\texttt{Header} ($n = 31$) and \texttt{Goal} ($n = 44$) are reported for completeness. Even
pooled over all ten matches their support is small enough that the figures should be read as
indicative rather than measured, and this is why both an unweighted and a support-weighted
mean are given: the unweighted mean assigns \texttt{Header} the same weight as \texttt{Pass},
which has $9{,}319$ instances.

\paragraph{A caveat that the per-match table makes visible.}
The value of the ball track depends on whether it is intact. In 128057 the ball leaves the
field of play on $3\%$ of frames, as it must; in 132831 and 132877 the recorded position is
clamped to the pitch rectangle and never does. On the Challenge fold, adding the ball raises
mean average precision at $\tau = 1\,\text{s}$ from $0.503$ to $0.775$ on 128057 and lowers it
from $0.369$ to $0.336$ on 132831. Any with-ball figure that pools a clean and a censored
match should be read with that decomposition rather than in isolation.
